% !TEX root = ../main.tex
% ESEMPIO 3: la stessa figura riusata in più frame.
%   La figura (examples/fig_reuse.tex) usa due coppie di numeri di overlay, First e Second. Ogni frame le
%   imposta come serve: la coppia che non deve comparire viene posta a 0, e siccome gli overlay partono da 1
%   non è mai visibile. Così un solo file di figura serve più slide, ognuna con la sua animazione.

\begin{frame}
    \frametitle{Example: one figure, first frame}

    \pgfmathsetmacro\numberFirstStart{int(2)}
    \pgfmathsetmacro\numberFirstEnd{int(\numberFirstStart+1)}
    % non usati in questo frame: a 0 non compaiono mai
    \pgfmathsetmacro\numberSecondStart{int(0)}
    \pgfmathsetmacro\numberSecondEnd{int(0)}

    \begin{adjustbox}{width=.8\linewidth, max height=.6\textheight, center}
        \begin{tikzpicture}[node distance=2.5em]
    \node[rboxes, on slide={<\numberFirstStart-\numberFirstEnd> highlight light}] (a) {Element A};
    \node[rboxes, right=of a, on slide={<\numberFirstStart-\numberFirstEnd> highlight light}] (b) {Element B};
    \node[rboxes, right=of b, on slide={<\numberSecondStart-\numberSecondEnd> highlight light}] (c) {Element C};

    \draw[->, link, on slide={<\numberFirstStart-\numberFirstEnd> highlight link}] (a) -- node[on link] {relation} (b);
    \draw[->, link, on slide={<\numberSecondStart-\numberSecondEnd> highlight link}] (b) -- node[on link] {relation} (c);
\end{tikzpicture}

    \end{adjustbox}

\end{frame}

\begin{frame}
    \frametitle{Example: the same figure, second frame}

    \pgfmathsetmacro\numberFirstStart{int(0)}
    \pgfmathsetmacro\numberFirstEnd{int(0)}
    \pgfmathsetmacro\numberSecondStart{int(2)}
    \pgfmathsetmacro\numberSecondEnd{int(\numberSecondStart+0)}

    \begin{adjustbox}{width=.8\linewidth, max height=.6\textheight, center}
        \begin{tikzpicture}[node distance=2.5em]
    \node[rboxes, on slide={<\numberFirstStart-\numberFirstEnd> highlight light}] (a) {Element A};
    \node[rboxes, right=of a, on slide={<\numberFirstStart-\numberFirstEnd> highlight light}] (b) {Element B};
    \node[rboxes, right=of b, on slide={<\numberSecondStart-\numberSecondEnd> highlight light}] (c) {Element C};

    \draw[->, link, on slide={<\numberFirstStart-\numberFirstEnd> highlight link}] (a) -- node[on link] {relation} (b);
    \draw[->, link, on slide={<\numberSecondStart-\numberSecondEnd> highlight link}] (b) -- node[on link] {relation} (c);
\end{tikzpicture}

    \end{adjustbox}

\end{frame}
